\newpage
\section{Discussion}
% ──────────────────────────────────────────────────────────────────────────────

The results reveal both encouraging convergences and persistent gaps in how
medical XAI is currently evaluated. We interpret the key findings below,
discuss limitations of the present study, and identify directions for future
work.

% ──────────────────────────────────────
\subsection{Interpretation of Key Findings}

Expanding the corpus to 52 studies introduces an important revision to the
Faithfulness finding. Rather than appearing universally, Faithfulness is
now present in 47 of 52 studies (90\%), with the five exceptions
concentrated in human-centred and consistency-focused
designs~\cite{rezaeian2026explainability,finzel2024human,bauer2025human,qin2025consistency,alpsalaz2026explainability}.
This pattern is analytically significant: Faithfulness absence is not
random but co-occurs with a shift in evaluation layer, signalling a genuine
conceptual divide between technical-fidelity and clinician-trust research
paradigms. Because Faithfulness now has non-trivial variance, Spearman
correlations involving it are computable. The F--S correlation is modest
but statistically significant ($\rho = 0.33$, $p = 0.016$), confirming
co-occurrence without implying full redundancy. The strongest pairwise
relationship in the expanded corpus is between AOPC/MoRF and Pixel
Flipping ($\rho = 0.51$, $p < 0.001$), which is mechanistically expected:
both are pixel-level perturbation methods that degrade saliency-ranked
features and measure the resulting prediction drop. AOPC/MoRF and IROF
also correlate significantly ($\rho = 0.37$, $p = 0.007$). If these
perturbation metrics reliably track the same underlying property,
reporting all three adds little information while increasing methodological
complexity. A minimal evaluation set anchored on Faithfulness, supplemented
by one perturbation-based metric such as AOPC, and one sensitivity measure,
may be sufficient for technical layer assessment in most medical imaging
contexts.

The isolation of Sufficiency remains equally telling. It appeared in only
4 of 52 studies (8\%) --- a proportion essentially unchanged from the
original corpus --- and exhibited negative correlations with all other
metrics ($\rho$ ranging from $-0.09$ to $-0.15$, all non-significant).
Its low co-occurrence is not simply a matter of adoption lag; it reflects
a conceptual gap. Sufficiency asks whether patients with similar clinical
profiles receive consistent explanations---a question of direct clinical
relevance---yet the literature has not treated it as a priority. This
disconnect echoes the broader critique by Kadir et al.~\cite{kadir2023evaluation}
that human-centered assessment remains underdeveloped relative to
quantitative metrics.

The sub-domain disparity between oncology and neurology further underscores
that metric richness is unevenly distributed. Where oncology benefits from
well-established imaging benchmarks that support diverse evaluation
protocols, neurology lacks analogous shared resources, which likely
suppresses both metric diversity and cross-study comparability.

% ──────────────────────────────────────
\subsection{Limitations}

Several limitations qualify these findings. First, the corpus of 52 studies,
while spanning all four sub-domains, is modest in absolute size and should be
treated as a focused scoping review rather than a definitive survey; the
patterns observed require replication on a larger systematically assembled
sample. Second, the co-occurrence analysis captures which metrics tend to
appear together but cannot establish whether two frequently paired metrics
are genuinely redundant or merely reported together by convention. Third, the taxonomy's three layers
were constructed deductively from prior literature, and while the corpus
mapping provides partial validation, a prospective user study with
clinical practitioners would be needed to confirm that the human-centered
trust layer captures what domain experts actually require from
explanations. Finally, restricting scope to medical XAI, while
analytically necessary, limits the generalizability of the findings
to other high-stakes domains. Work such as Ghosh et
al.~\cite{Ghosh2024DemystifyingEC}, who apply novel XAI metrics to
industrial IC packaging inspection, and Nagy et al.~\cite{Nagy2024ImprovingUM},
who use XAI metrics to improve accuracy in urban mapping surveys, suggests
that the metric redundancy and coverage gaps we document in medicine may
be symptomatic of a broader cross-domain problem worth investigating.

\subsection{Future Work}

Three directions follow directly from the present findings. First,
a prospective study instrumenting clinicians as they interact with
SHAP- and Grad-CAM-generated explanations would ground the trust layer
of the taxonomy in behavioral evidence. Second, developing a standardized
reporting template for medical XAI evaluations---specifying a minimal
metric set alongside optional domain-specific additions---could reduce
the fragmentation this review documents. Third, extending future analyses
to include user-study metrics such as cognitive load and decision accuracy
would allow researchers to test empirically whether high Faithfulness
scores translate into genuine clinical utility, a question the current
literature leaves largely unanswered.
\section{Conclusion}
% ──────────────────────────────────────────────────────────────────────────────

This paper investigated the evaluation landscape of XAI within the medical
domain, focusing on how quantitative metrics co-occur and distribute
across oncology, cardiology, neurology, and radiology. Drawing on a
focused corpus of 35 peer-reviewed studies, we identified a dominant cluster of
technically-grounded metrics---Faithfulness and Sensitivity---that
consistently co-occur and collectively account for the majority of evaluation
activity. Perturbation-based metrics such as AOPC and IROF form a secondary
cluster, while Sufficiency remains an isolated and underutilised criterion
despite its direct clinical relevance.

Mapping the corpus onto a three-layer taxonomy of mathematical faithfulness,
system robustness, and human-centered trust confirmed that current evaluation
practice is heavily concentrated in the first layer. 47 of 52 studies addressed
mathematical faithfulness, but only 34\% addressed system robustness and
9\% included any human-centered criterion,
leaving significant gaps precisely in the dimensions most consequential
for clinical deployment.

Taken together, these findings support two main conclusions. First, a
minimal but comprehensive evaluation set---combining one faithfulness
measure, one perturbation-based measure, one sufficiency check, and at
least one human-centered criterion---would cover all three taxonomy layers
without imposing an unreasonable methodological burden. Second, the field
needs shared benchmarks and reporting standards that make cross-study
comparison tractable, particularly for underserved sub-domains such as
neurology where evaluation practice currently lags. Addressing these gaps
is a necessary step toward medical AI systems whose explanations are not
only technically faithful but genuinely trustworthy in clinical practice.

\newpage
% ─────────────────────────────────────────────────────────────────────────────
